\section{Discussion}
\label{sec:discussion}

The CIFAR-10 results are overwhelmingly negative for the variance-reduction composition hypothesis. This section analyzes each failure mode, explains the 2D-to-CIFAR gap, and answers our research questions.

\subsection{Why Cell\_000 (Baseline) Failed}

Cell\_000 uses the simplest configuration: independent coupling, standard CFM loss, and uniform timestep sampling. It diverged to NaN at step ${\sim}18{,}700$. The root cause is a numerical issue in the OT path's conditional vector field (Eq.~\ref{eq:ot_vf}):
\begin{equation}
    \ut(x | x_1) = \frac{x_1 - (1 - \sigma_{\min}) x}{1 - (1 - \sigma_{\min}) t}.
\end{equation}
The denominator $1 - (1 - \sigma_{\min}) t$ approaches $\sigma_{\min} = 10^{-4}$ as $t \to 1$. Under fp16 mixed precision, this produces values near $10^{4}$, which cascade through the network gradients and eventually overflow.

\textbf{Why BatchOT prevents this:} With independent coupling, the numerator $x_1 - (1-\sigma_{\min})x$ can be large because $x_0$ and $x_1$ are unrelated---a noise sample may be far from its paired data point. BatchOT coupling minimizes the transport cost, keeping $\|x_1 - x_0\|$ small. Since $x_t = t x_1 + \sigma_t x_0$ interpolates between paired points, both the numerator and denominator of $\ut$ remain moderate, preventing the overflow. This explains why cell\_100 converges while cell\_000 does not: BatchOT is not just a variance-reduction technique---it is a \emph{numerical stabilizer} for the OT path under limited precision.

\subsection{Why StableVM Failed}

All four StableVM cells (010, 011, 110, 111) exhibit extreme loss oscillation between ${\sim}40$ and ${\sim}22{,}000$. The cause lies in the log-probability computation (Eq.~\ref{eq:log_prob}):
\begin{equation}
    \log p_t(x | x_1) = -\frac{d}{2}\log(2\pi) - d \log \sigma_t - \frac{1}{2\sigma_t^2}\|x - \mu_t(x_1)\|^2.
\end{equation}

The critical term is $-d \log \sigma_t$. For CIFAR-10, $d = 3 \times 32 \times 32 = 3072$. Near $t = 1$, $\sigma_t \approx \sigma_{\min} = 10^{-4}$, so:
\begin{equation}
    -d \log \sigma_t \approx -3072 \times \log(10^{-4}) \approx 3072 \times 9.21 \approx 28{,}300.
\end{equation}
This enormous value dominates the softmax in the importance weights (Eq.~\ref{eq:stablevm_weights}), causing the weights to concentrate on a single reference sample---whichever $x_1^{(k)}$ is closest to $\xt$. The resulting target is effectively a single-sample estimate with high variance, but computed through a numerically unstable pathway. The loss oscillation reflects the softmax flipping between different reference samples across training steps.

\textbf{Why this is invisible in 2D:} When $d = 2$, the same term yields $-2 \times \log(10^{-4}) \approx 18.4$, which is small enough that the softmax produces meaningful weight distributions across all $K$ references. The technique works as intended in low dimensions but breaks catastrophically when $d$ is large.

\subsection{Why TPC Failed}

Cells 001 and 101 diverged to NaN at steps ${\sim}16{,}600$ and ${\sim}14{,}400$ respectively. The TPC consistency penalty (Eq.~\ref{eq:tpc_loss}) compares velocity predictions at complementary timesteps:
\begin{equation}
    \lambda \left\| \vt(\xt^{(1)}, t_1) - \vt(\xt^{(2)}, t_2) \right\|^2, \quad t_2 = 1 - t_1.
\end{equation}

When $t_1$ is close to 0 and $t_2$ close to 1, the velocity magnitudes at these two timesteps are fundamentally different: near $t=0$ the flow must move slowly from noise, while near $t=1$ it must make fine adjustments toward the data manifold. The conditional vector field magnitudes differ by orders of magnitude due to the $1/(1-(1-\sigma_{\min})t)$ denominator. The consistency penalty produces gradients that are proportional to this mismatch, leading to gradient explosions that destabilize training.

\textbf{Note:} Cell\_000 also diverged to NaN without TPC, but for a different reason (fp16 overflow from independent coupling). The NaN onset is earlier for TPC cells (14--17k steps vs.\ 18.7k), suggesting that TPC accelerates the instability.

\subsection{The 2D-to-CIFAR Gap}

All four flow-matching variants in Study~1 work perfectly in 2D, with no sign of the catastrophic failures observed on CIFAR-10. The fundamental reason is \textbf{dimensionality}:

\begin{itemize}[leftmargin=*]
    \item StableVM's log-probability scales as $O(d)$: harmless when $d=2$, catastrophic when $d=3072$.
    \item The OT path's denominator creates velocities of magnitude $\sim 1/\sigma_{\min}$. In 2D, the gradient norm stays bounded; in $d=3072$, the per-channel gradient accumulation amplifies the instability.
    \item TPC's consistency penalty couples velocities at $t$ and $1-t$ whose magnitudes diverge more severely in high dimensions, where the velocity field must navigate a higher-dimensional landscape.
\end{itemize}

This gap carries a methodological warning: \textbf{2D experiments are valuable for building geometric intuition but insufficient for validating numerical stability}. Any technique involving log-probabilities, importance weights, or time-endpoint comparisons should be stress-tested in the target dimension before deployment.

\subsection{Answering the Research Questions}

\paragraph{RQ1: Are the three axes additive?} \textbf{No.} The results are emphatically negative. Not only are the axes not additive, but combining BatchOT with either StableVM or TPC actively \emph{destroys} the convergence that BatchOT alone provides. Cell\_100 (BatchOT only) converges to FID~18.4; adding TPC (cell\_101) causes NaN divergence; adding StableVM (cell\_110) causes wild instability.

\paragraph{RQ2: Which single axis provides the largest improvement?} \textbf{BatchOT coupling}, unambiguously. It is the only axis whose inclusion leads to convergence (when used alone). StableVM and TPC both degrade performance in every tested configuration. BatchOT acts as both a variance reducer (shorter, straighter trajectories) and a numerical stabilizer (bounded velocity magnitudes).

\paragraph{RQ3: Do 2D insights transfer?} \textbf{Partially.} The geometric insight from Study~1---that OT-type paths produce straighter trajectories and better low-NFE robustness---transfers directly: cell\_100's midpoint solver at NFE=40 matches Euler at NFE=100, confirming the trajectory-straightness advantage. However, the stability insights do \emph{not} transfer: techniques that work flawlessly in 2D fail catastrophically in high dimensions.
